\documentclass[12pt, ukrainian]{article}
\usepackage[a4paper]{geometry}
\setlength\parindent{0mm}
\geometry{top=1.1cm,bottom=1.1cm,left=1.1cm,right=1.1cm}
\pagestyle{empty}
\usepackage[T2A]{fontenc}
\usepackage[utf8]{inputenc}
\usepackage[ukrainian]{babel}
\usepackage{graphicx}
\usepackage{caption}
\usepackage{verbatim, listings, clrscode3e, fancyvrb}
\DeclareCaptionFormat{nocaption}{}
\captionsetup{format=nocaption,aboveskip=0pt,belowskip=0pt}
\usepackage[Export]{adjustbox}
\adjustboxset{max size={0.9\linewidth}{0.9\paperheight}}
\def\code#1{\Verb!#1!}
\begin{document}
{{16.05.2024 Контрольна робота, варіант  \No \textbf { 18 }} група К\_\_ ПІБ {\parbox {2.5 cm} \dotfill}\par}
{
%\begin {large}
\begin{minipage}[t]{9.7cm}
У завданнях 1--3 вказати, що буде виведено за виконання. Повідомлення інтерпретатора про помилку позначати ERROR


               
{\begin{center}Завдання {1}\end{center}}
\vspace{-4mm}

\begin{verbatim}
e = ('a','b','c','d','e',0,1)
print(e[-2])

\end{verbatim}            

\vspace{15mm}
               
{\begin{center}Завдання {2}\end{center}}
\vspace{-4mm}

\begin{verbatim}
c=['0','1','2','3']
c[:-8]=[1,0]
print(c)

\end{verbatim}            


\end{minipage}
\vline \hspace{8mm}
\begin{minipage}[t]{8.0cm}                
               
{\begin{center}Завдання {3}\end{center}}
\vspace{-4mm}

\begin{verbatim}
_d = 0

class B:
    _d = 1
    
    def __init__(self):
        self._d = 2
        _d = 3
        B._d = 4

obj = B()
obj._d = 5
try:
    print(_d, end=' ')
    print(obj._d, end=' ')
    print(B._d, end=' ')
    print(obj._d, end=' ')
    print(obj._B__d)
except:
    print('ERR')

\end{verbatim}            

\end{minipage}


               
{\begin{center}Завдання {4}\end{center}}
\vspace{-4mm}
Нехай змінна \kw{s} позначає дуже довгу послідовність цілих, по якій можна ітерувати.

Написати вираз-генератор, що задає підпослідовність її елементів, що не діляться на жодне з чисел 2, 7, 5.

Обмеження: усі елементи шуканої підпослідовності одночасно в пам'яті розміщені бути не можуть.\par

\vspace{10mm}
               
{\begin{center}Завдання {5}\end{center}}
\vspace{-4mm}
Змінні \kw{next} та \kw{prev} вузла списку позначають наступний та попередній за ним вузол відповідно. Починаючи з голови, у список записано послідовні цілі числа від~$-51$ до~$35$. Нехай змінна \kw{p} позначає вузол, що зберігає значення~$-5$.
Яке число зберігається у вузлі \kw{p.next.next.prev.prev.next} ?\par

               
{\begin{center}Завдання {6}\end{center}}
\vspace{-4mm}
Значенням змінної \kw{a} є цілочисловий масив типу $numpy.ndarray$ розмірів $4{\times}7$. На скільки збільшиться сума елементів масиву \kw{a} після виконання
\begin{verbatim}
a[-2] += 1
\end{verbatim}


Якщо код некоректний, то в якості відповіді зазначити ERROR.\par

               
{\begin{center}Завдання {10}\end{center}}
\vspace{-4mm}

Написати функцію $f$, яка для файлу, заданого шляхом, розв’язує задачу:\\ Текстовий файл містить послідовність чисел $a_1$, $a_2$, $\ldots$, $a_t$, записану у форматі одне число на рядок. Знайти третє найменше значення, що записане у файлі. Кожне значення, що записане у файлі, враховується стільки разів, скільки входить у файл. Для послідовності 2, 4, 12, 5, 2, 0 таким значенням буде 2.

В усіх випадках розміри файлів такі, що їх вміст повністю завантажений у пам'ять бути не може. Якість коду також оцінюється.



\vspace{10mm}\par

\newpage
\begin{minipage}[t]{9.7cm}
               
{\begin{center}Завдання {7}\end{center}}
\vspace{-4mm}

\begin{center}
	\adjustimage{max size={0.8\linewidth}{0.8\paperheight}}
	{../15BinTree/trees_exp/77.png}
\end{center}
Указати послідовність міток вершин дерева, зображеного на рис., що утвориться в результаті обходу дерева в прямому порядку. Мітки записувати через пробіл.\par Алгоритм починає роботу з кореня (на рис. це верхня вершина).\par

\end{minipage}
\vline \hspace{8mm}
\begin{minipage}[t]{8.0cm}
               
{\begin{center}Завдання {8}\end{center}}
\vspace{-4mm}
Для графа на рисунку з вершини 3 запущено алгоритм пошуку в ширину, що будує кістякове дерево. Списки суміжності вершин переглядаються алгоритмом за зростанням міток вершин.\par
\begin{center}
	\adjustimage{max size={0.9\linewidth}{0.9\paperheight}}
	{../16Graphs/Traversals/348.png}
\end{center}
\par Які з наведених ребер входять до складу отриманого кістякового дерева?\par
1) (0, 3)\par
2) (0, 5)\par
3) (3, 5)\par
4) (3, 4)\par
5) (2, 3)\par
6) (1, 6)\par
{\vskip 15pt} \par

\end{minipage}

               
{\begin{center}Завдання {9}\end{center}}
\vspace{-4mm}
Рядки журналу через кому містять ім'я користувача (складається з латиниці), час події (фор\-мат як у прикладі), тип події (ERROR, WARNING, INFO, STATUS), код події (ціле).

Білі символи навколо роздільників не допускаються.

Приклад такого рядка присвоєно змінній \kw{s}.



\begin{verbatim}
s = 'username,2022-05-05 13:26,ERROR,404'
\end{verbatim}


Нижче наведено код, за виконання якого значенням змінної \kw{out} має стати код події.

Для наведеного прикладу значенням змінної \kw{out} має стати рядок \code{404}



\begin{verbatim}
res = re.fullmatch(r'XXX', s)
s = ''
out = YYY
\end{verbatim}


Код має працювати не тільки для конкретного рядка з прикладу, але й для кожного рядка з журналу.


Який рядок треба записати на місце \kw{XXX} ?
Який рядок треба записати на місце \kw{YYY} ?\par

%\end {large}
}
%end
\end{document}

